\hypertarget{cryptography-and-hashing-hashlib-hmac}{%
\chapter{\texorpdfstring{Cryptography and Hashing (\texttt{hashlib},
\texttt{hmac})}{Cryptography and Hashing (hashlib, hmac)}}\label{cryptography-and-hashing-hashlib-hmac}}

Security-sensitive hashing and Message Authentication Codes (MACs) are
handled by \passthrough{\lstinline!hashlib!} and
\passthrough{\lstinline!hmac!}, which act as bridges to the system's
OpenSSL library.

\hypertarget{hashlib-the-openssl-bridge}{%
\subsection{\texorpdfstring{45.1 \texttt{hashlib}: The OpenSSL
Bridge}{45.1 hashlib: The OpenSSL Bridge}}\label{hashlib-the-openssl-bridge}}

\passthrough{\lstinline!hashlib!} provides a common interface to many
different secure hash and message digest algorithms.

\hypertarget{static-vs.-dynamic-algorithms}{%
\subsubsection{1. Static vs.~Dynamic
Algorithms}\label{static-vs.-dynamic-algorithms}}

\begin{itemize}
\tightlist
\item
  \textbf{Guaranteed}: \passthrough{\lstinline!sha256!},
  \passthrough{\lstinline!sha512!}, \passthrough{\lstinline!md5!} are
  always available.
\item
  \textbf{OpenSSL-dependent}: Algorithms like
  \passthrough{\lstinline!blake2b!} or \passthrough{\lstinline!sha3!}
  are available only if the linked OpenSSL library supports them.
\end{itemize}

\hypertarget{releasing-the-gil}{%
\subsubsection{2. Releasing the GIL}\label{releasing-the-gil}}

Hashing large files can be CPU-intensive. CPython's
\passthrough{\lstinline!hashlib!} implementations \textbf{release the
GIL} during the \passthrough{\lstinline!update()!} call if the data is
large enough. This allows true parallelism when hashing multiple files
in separate threads.

\hypertarget{hmac-keyed-hashing-for-message-authentication}{%
\subsection{\texorpdfstring{45.2 \texttt{hmac}: Keyed-Hashing for
Message
Authentication}{45.2 hmac: Keyed-Hashing for Message Authentication}}\label{hmac-keyed-hashing-for-message-authentication}}

\passthrough{\lstinline!hmac!} implements the HMAC algorithm as defined
by RFC 2104. * \textbf{Why not just
\passthrough{\lstinline!hash(key + message)!}?}: Simple concatenation is
vulnerable to ``length-extension attacks'' in certain hash functions
(like MD5 and SHA-1). \passthrough{\lstinline!hmac!} uses a
double-hashing nested structure to prevent this. *
\textbf{\passthrough{\lstinline!compare\_digest(a, b)!}}: Always use
this function for comparing hashes/tokens. It is a
\textbf{constant-time} comparison, preventing ``timing attacks'' where
an attacker can deduce the correct token by measuring how long the
comparison takes to fail.

\begin{center}\rule{0.5\linewidth}{0.5pt}\end{center}
